\documentclass[11pt]{article}
\usepackage{fontspec}
\setmainfont{Times New Roman}
\setsansfont{Arial}
\setmonofont{Consolas}
\usepackage{amsmath,amssymb,mathtools}
\usepackage{geometry}
\usepackage{microtype}
\newcommand{\mo}{\mathfrak{o}}
\newcommand{\mO}{\mathfrak{O}}
\newcommand{\mT}{\mathfrak{T}}
\newcommand{\mR}{\mathfrak{R}}
\newcommand{\mq}{\mathfrak{q}}
\geometry{a4paper,margin=2.35cm}
\setlength{\parindent}{1.25em}
\setlength{\parskip}{0pt}
\makeatletter
\renewcommand\footnotesize{\@setfontsize\footnotesize{7.7}{8.7}\selectfont}
\makeatother
\emergencystretch=100em
\allowdisplaybreaks
\sloppy
\hbadness=10000
\vbadness=10000
\hfuzz=2pt
\vfuzz=10pt
\raggedbottom

\begin{document}

% Unit: Paper32 Schur paragraph v001. Direct Russian translation from clean RA34 Paper32 line 8 / segment P32-S0004, checked against source scan page 1 and Claude batch OCR witness.
\setcounter{footnote}{0}

\noindent Вопрос о числовых полях наименьшей степени над основным полем $P$, в которых неприводимое над $P$ представление конечной группы распадается на абсолютно неприводимые составляющие, впервые рассмотрел И. Шур.\footnote{\emph{Arithmetische Untersuchungen über endliche Gruppen}, Sitzungsber. d. Berl. Ak. d. Wiss. 1906, с. 164. \emph{Beiträge zur Theorie der Gruppen linearer Substitutionen}, Transact. of the Am. Math. Soc. Ser. 2, Vol. XV, 1909, с. 159. Во второй работе результаты перенесены на вполне приводимые гиперкомплексные системы. В один класс объединяется представление со всеми его преобразованными (подобными).} Пусть, например, $P$ содержит характер такой составляющей. Тогда И. Шур показал, что степень этих полей над $P$ -- индекс -- совпадает с числом тех абсолютно неприводимых составляющих, которые все принадлежат одному и тому же классу; что, следовательно, то же самое поле уже определяется требованием выделить одну абсолютно неприводимую составляющую; далее, что степень всякого поля над $P$, в котором происходит такой распад, является кратной индексу. Все эти поля будем называть полями расщепления, также и в случае, когда $P$ не содержит характера. На примере И. Шур показал, что поля расщепления наименьшей степени не обязаны быть изоморфными. Если $P$ не содержало соответствующего неприводимого характера, то поля расщепления должны содержать поле такого характера и его сопряженные относительно $P$; абсолютно неприводимые представления, принадлежащие различным сопряженным характерам, хотя и сопряжены, но, очевидно, принадлежат различным классам. Для выделения системы абсолютно неприводимых представлений одного класса и здесь достаточно поля соответствующего характера и одного из соответствующих полей расщепления.

\end{document}
